%============================================================
%  Bd. 24 - Schlichte ungerichtete Graphen
%============================================================

\documentclass[main.tex]{subfiles}

\ifSubfilesClassLoaded{
  \BandSetupStandalone{B24}
}{
}

\title{Bd. 24 - Schlichte ungerichtete Graphen}
\author{Martin Kunze}
\date{}
\setcounter{file}{24}

\hypersetup{
  pdftitle={Bd. 24 - Schlichte ungerichtete Graphen},
  pdfauthor={Martin Kunze}
}

\begin{document}

\title{Bd. 24 - Schlichte ungerichtete Graphen}
\author{Martin Kunze}
\date{}
\hypersetup{pageanchor=false}
\maketitle
\hypersetup{pageanchor=true}
\setcounter{tocdepth}{3}
\tableofcontents
\setcounter{file}{24}
\renewcommand{\FormulaBandID}{24}

\chapter{Einleitung}

Ein schlichter ungerichteter Graph ist ein ungerichteter Graph ohne
Knotenschleifen. Die Schleifenfreiheit wird als eigenes Strukturaxiom
geführt, sodass spätere Beweise sowohl auf die Symmetrie des Grundgraphen
als auch unmittelbar auf die Irreflexivität seiner Kantenrelation zugreifen
können.

Der Band zeigt außerdem, dass induzierte Teilgraphen schlichter Graphen
wieder schlicht sind. Damit steht die in der Baumtheorie benötigte
Graphklasse unabhängig vom Zusammenhangsbegriff des folgenden Bandes zur
Verfügung.

\chapter{Schlichte ungerichtete Graphen}

\section{Axiomatischer Begriff}

\FormulaDefDeltaK[Begriff des schlichten ungerichteten Graphen]{%
  \SimpleUndirectedGraphStruct{V,E}%
}{SimpleUndirectedGraphDef}{%
  \DeltaRow{Mengen}{V\dsep E\dsep x}
  \DeltaRow{\textbf{Axiome}}{}
  \DeltaRow{Ungerichteter Graph}{\UndirectedGraphStruct{V,E}}
    [\FormulaRefAuto{UndirectedGraphDef}[def]]
  \DeltaRow{Schleifenfreiheit}{x\in V\vdash(x,x)\notin E}
  \DeltaRow{\textbf{Neues Symbol}}{}
  \DeltaPrem{Schlichter ungerichteter Graph}{%
    \SimpleUndirectedGraphStruct{V,E}}
}

\FormulaAxiomDeltaK[Zugriff auf den ungerichteten Grundgraphen]{%
  \SimpleUndirectedGraphStruct{V,E}\vdash
  \UndirectedGraphStruct{V,E}%
}{SimpleUndirectedGraphUnderlyingAxiom}{%
  \DeltaRow{Mengen}{V\dsep E}%
}

\FormulaAxiomDeltaK[Zugriff auf die Schleifenfreiheit]{%
  \SimpleUndirectedGraphStruct{V,E}\vdash
  \forall x\in V\,\bigl((x,x)\notin E\bigr)%
}{SimpleUndirectedGraphIrreflexiveAxiom}{%
  \DeltaRow{Mengen}{V\dsep E\dsep x}%
}

\FormulaThmDeltaK[Strukturzugriff auf einen schlichten Graphen]{%
  \SimpleUndirectedGraphStruct{V,E}\vdash
  \UndirectedGraphStruct{V,E}\land\GraphStruct{V,E}%
}{SimpleUndirectedGraphAccess}{%
  \DeltaRow{Mengen}{V\dsep E}%
}
\begin{tabproof}
  \proofstep{1}{\SimpleUndirectedGraphStruct{V,E}}{\rA}
  \proofstep{1}{\UndirectedGraphStruct{V,E}}{%
    \FormulaRefAuto{SimpleUndirectedGraphUnderlyingAxiom}[ax]{1}}
  \proofstep{1}{\GraphStruct{V,E}}{%
    \FormulaRefAuto{UndirectedGraphAccess}[thm]{2}}
  \proofstep{1}{\UndirectedGraphStruct{V,E}\land\GraphStruct{V,E}}{\rAI{2,3}}
\end{tabproof}

\FormulaThmDeltaK[Einführungskriterium für schlichte ungerichtete Graphen]{%
  \begin{aligned}[t]
    &E\subseteq V\times V\dsep
      \forall x,y\in V\,\bigl((x,y)\in E\rightarrow(y,x)\in E\bigr)\dsep{}\\[-2pt]
    &\forall x\in V\,\bigl((x,x)\notin E\bigr)
      \vdash\SimpleUndirectedGraphStruct{V,E}
  \end{aligned}%
}{SimpleUndirectedGraphIntroduction}{%
  \DeltaRow{Mengen}{V\dsep E\dsep x\dsep y}%
}
\begin{tabproofwide}
  \proofstepwidestar[1]{E\subseteq V\times V}{\rA}
  \proofstepwidestar[2]{%
    \forall x,y\in V\,\bigl((x,y)\in E\rightarrow(y,x)\in E\bigr)}{\rA}
  \proofstepwidestar[3]{\forall x\in V\,\bigl((x,x)\notin E\bigr)}{\rA}
  \proofstepwidestar[1,2]{\UndirectedGraphStruct{V,E}}{%
    \FormulaRefAuto{UndirectedGraphIntroduction}[thm]{1,2}}
  \proofstepwidestar[5]{x\in V}{\rA}
  \proofstepwidestar[3,5]{(x,x)\notin E}{\rRE{5,3}}
  \proofstepwidestar[3]{\forall x\in V\,\bigl((x,x)\notin E\bigr)}{%
    \rUI{\rRI{5,6}}}
  \proofstepwidestar[1,2,3]{\SimpleUndirectedGraphStruct{V,E}}{%
    \FormulaRefAuto{SimpleUndirectedGraphDef}[def]{4,7}}
\end{tabproofwide}

\FormulaThmDeltaK[Ein schlichter Graph besitzt keine Knotenschleife]{%
  \SimpleUndirectedGraphStruct{V,E}\dsep x\in V
  \vdash(x,x)\notin E%
}{SimpleUndirectedGraphLoopFree}{%
  \DeltaRow{Mengen}{V\dsep E\dsep x}%
}
\begin{tabproof}
  \proofstep{1}{\SimpleUndirectedGraphStruct{V,E}}{\rA}
  \proofstep{2}{x\in V}{\rA}
  \proofstep{1}{\forall x\in V\,\bigl((x,x)\notin E\bigr)}{%
    \FormulaRefAuto{SimpleUndirectedGraphIrreflexiveAxiom}[ax]{1}}
  \proofstep{1,2}{(x,x)\notin E}{%
    \rRE{2,\rUE{3}}}
\end{tabproof}

\section{Induzierte Teilgraphen}

\FormulaThmDeltaK[Induzierte Teilgraphen schlichter Graphen sind schlicht]{%
  \SimpleUndirectedGraphStruct{V,E}\dsep U\subseteq V\vdash
  \SimpleUndirectedGraphStruct{U,\DirectedInducedEdges{E}{U}}%
}{SimpleUndirectedInducedSubgraph}{%
  \DeltaRow{Mengen}{V\dsep E\dsep U\dsep x}%
}
\begin{tabproofwide}
  \proofstepwidestar[1]{\SimpleUndirectedGraphStruct{V,E}}{\rA}
  \proofstepwidestar[2]{U\subseteq V}{\rA}
  \proofstepwidestar[1]{\UndirectedGraphStruct{V,E}}{%
    \FormulaRefAuto{SimpleUndirectedGraphUnderlyingAxiom}[ax]{1}}
  \proofstepwidestar[1,2]{%
    \UndirectedGraphStruct{U,\DirectedInducedEdges{E}{U}}}{%
    \FormulaRefAuto{UndirectedInducedSubgraph}[thm]{3,2}}
  \proofstepwidestar[5]{x\in U}{\rA}
  \proofstepwidestar[2,5]{x\in V}{%
    \FormulaRefAuto{A\subseteq B\dsep x\in A\vdash x\in B}[thm]{2,5}}
  \proofstepwidestar[1,2,5]{(x,x)\notin E}{%
    \FormulaRefAuto{SimpleUndirectedGraphLoopFree}[thm]{1,6}}
  \proofstepwidestar[8]{(x,x)\in\DirectedInducedEdges{E}{U}}{\rA}
  \proofstepwidestar[1,2,8]{%
    (x,x)\in E\land x\in U\land x\in U}{%
    \FormulaRefAuto{P\leftrightarrow Q, P\vdash Q}[thm]{%
      \FormulaRefAuto{DirectedInducedEdgeMembership}[thm]{%
        \FormulaRefAuto{UndirectedGraphAccess}[thm]{3},2},8}}
  \proofstepwidestar[1,2,8]{(x,x)\in E}{\rAEa{9}}
  \proofstepwidestar[1,2,5,8]{\bot}{\rBI{10,7}}
  \proofstepwidestar[1,2,5]{%
    (x,x)\notin\DirectedInducedEdges{E}{U}}{\rCI{8,11}}
  \proofstepwidestar[1,2]{%
    \SimpleUndirectedGraphStruct{U,\DirectedInducedEdges{E}{U}}}{%
    \FormulaRefAuto{SimpleUndirectedGraphDef}[def]{4,12}}
\end{tabproofwide}

\end{document}
